% wave12_a1_bar_cobar_gelfand.tex
% Adversarial assault and rigorous reconstruction of the bar--cobar
% foundations and CY-A_3 for Vol~III. Five attack<->heal cycles.
% Channels: I.~M.~Gelfand, directness; Beilinson, dismissal of the false;
% Costello--Francis--Gwilliam, factorisation-algebra primary source.
% Companion: working_notes.tex Sections ``coderived status'',
% cy_to_chiral.tex Theorems thm:cy-to-chiral, thm:cy-to-chiral-d3,
% thm:derived-framing-obstruction, Prop:bar-ce-chiral.

\providecommand{\ClaimStatusTheorem}{\textsuperscript{\textsc{[th]}}}
\providecommand{\ClaimStatusConjectured}{\textsuperscript{\textsc{[cj]}}}
\providecommand{\ClaimStatusDerivation}{\textsuperscript{\textsc{[dv]}}}
\providecommand{\ClaimStatusFalsified}{\textsuperscript{\textsc{[fx]}}}
\providecommand{\kch}{\kappa_{\mathrm{ch}}}
\providecommand{\kcat}{\kappa_{\mathrm{cat}}}
\providecommand{\Ran}{\mathrm{Ran}}
\providecommand{\HH}{\mathrm{HH}}
\providecommand{\CE}{\mathrm{CE}}
\providecommand{\Conv}{\mathrm{Conv}}
\providecommand{\MC}{\mathrm{MC}}
\providecommand{\Einf}{E_\infty}
\providecommand{\Eone}{E_1}
\providecommand{\Etwo}{E_2}
\providecommand{\Ethree}{E_3}
\providecommand{\Ainf}{A_\infty}
\providecommand{\Linf}{L_\infty}
\providecommand{\Sym}{\mathrm{Sym}}
\providecommand{\Coh}{\mathrm{Coh}}
\providecommand{\Perf}{\mathrm{Perf}}
\providecommand{\C}{\mathbb{C}}
\providecommand{\Z}{\mathbb{Z}}
\providecommand{\Q}{\mathbb{Q}}
\providecommand{\RHom}{\mathrm{RHom}}
\providecommand{\Rep}{\mathrm{Rep}}
\providecommand{\cA}{\mathcal{A}}
\providecommand{\cC}{\mathcal{C}}
\providecommand{\cH}{\mathcal{H}}
\providecommand{\cM}{\mathcal{M}}
\providecommand{\cO}{\mathcal{O}}
\providecommand{\cL}{\mathcal{L}}
\providecommand{\fL}{\mathfrak{L}}
\providecommand{\fg}{\mathfrak{g}}
\providecommand{\fh}{\mathfrak{h}}
\providecommand{\fgl}{\mathfrak{gl}}
\providecommand{\Alg}{\mathrm{Alg}}
\providecommand{\barB}{\overline{B}}
\providecommand{\barA}{\overline{A}}
\providecommand{\Zder}{Z^{\mathrm{der}}_{\mathrm{ch}}}
\providecommand{\Dg}[1]{D_{\mathrm{g=#1}}}
\providecommand{\Heis}{\mathrm{Heis}}

\section*{Adversarial reconstruction of $\Omega \circ B = \mathrm{id}$, of
  $B(U^{\mathrm{ch}}(\fL)) \simeq \CE_*(\fL)$, and of CY-A$_3$}

\subsection*{Gelfand's discipline}

A theorem is one identity. Either it holds on the nose, or it
holds up to a controlled homotopy, or it does not hold. A statement
that has not been reduced to the first two shapes is not yet a
theorem. We work through the bar--cobar foundations of Vol~III by
asking, of each equality sign, which of the three it is, and under
what precise hypothesis on the input. The exercise falsifies three
standing formulations, sharpens four proofs, and exposes one
irreducible gap.

% ============================================================
\section*{Cycle 1. $\Omega(B(A)) \simeq A$ --- when, on the nose?}
% ============================================================

\paragraph{Attack.} Vol~I Theorem~A constructs an adjunction
$\Omega \dashv B$ between chiral algebras and coaugmented conilpotent
chiral coalgebras on $\Ran(X)$; Vol~I Theorem~B asserts the counit
$\psi \colon \Omega(B(\cA)) \xrightarrow{\sim} \cA$ is a
quasi-isomorphism \emph{on the Koszul locus}. Vol~III invokes the
inversion identity repeatedly in the CY-to-chiral construction
(Theorem~\ref{thm:cy-to-chiral-d3}; Proposition~\ref{prop:en-envelope-preservation}).
The hypothesis ``Koszul locus'' is not an empty restriction:
it is a \emph{genus-zero bar-cohomology concentration} statement.
The adversarial question is: does the inversion hold for
the chiral avatar of a Calabi--Yau origin, and in which dimension?

Falsification attempt. Take $\cA = \cH_{\mathrm{Muk}}$, the rank-$24$
Mukai Heisenberg VOA over a K3 surface.  Over a genus-$0$ curve, the
bar complex $B^{(0)}(\cA)$ has $m_0^{(0)} = 0$ and the cobar counit
is a qi by Vol~I, Theorem~B. Over a genus-$g$ curve,
$m_0^{(g)} = \kch \cdot \omega_g$ with $\kch = 2$, so the bar
differential acquires a nonzero curvature. The inversion
identity then moves from the category of chiral dg-coalgebras on
$\Ran$ to Positselski's coderived category
$D^{\mathrm{co}}(A)$. In this larger category, the
adjunction \emph{survives}; in the smaller derived category, it
does not. The abstract curved-coderived Quillen equivalence of
Booth--Lazarev (arXiv:2304.08409) is a theorem of associative
chain-complexes; it does not automatically descend to the
Ran-space / factorisation-structure setting.

\paragraph{Heal.}

\begin{theorem}[Scope of chiral inversion, genus-stratified]
\label{thm:gelfand-inversion-genus}
\ClaimStatusDerivation{}
Let $\cA$ be an augmented chiral algebra on a smooth complex curve
$C$ of genus $g$, with modular characteristic $\kch = \kch(\cA)$.
\begin{enumerate}[label=\textup{(\roman*)}]
\item\label{inv:g0} At $g = 0$: if $\cA$ is Koszul
(i.e.\ $H^{\geq 2} \barB(\cA) = 0$ over the genus-$0$ curve), the
counit
$\psi \colon \Omega(B(\cA)) \to \cA$
is a quasi-isomorphism in the ordinary derived category of chiral
algebras on $\Ran(C)$. The identity is on the nose at the
$\infty$-categorical level and up to an explicit Koszul-inversion
homotopy at the chain level.

\item\label{inv:ghigher} At $g \geq 1$ with $\kch \neq 0$: the
genus-$g$ bar complex $\barB^{(g)}(\cA)$ is \emph{curved} with
curvature $m_0^{(g)} = \kch \cdot \omega_g \neq 0$. The
ordinary derived category of chiral dg-coalgebras on $\Ran(C)$
does not house this datum. The inversion identity is
well-formed only in a coderived category $D^{\mathrm{co},g}(\cA)$
in the sense of Positselski.

\item\label{inv:chiral-bl} The extension of the Booth--Lazarev
curved Quillen equivalence (associative chain-complexes) to
chiral dg-coalgebras on $\Ran(C)$ is \emph{conditional} on
three chiral instantiation steps (Sections C1--C2 of
\texttt{working\_notes.tex} Subsection~\ref{subsec:coderived-status}):
coaugmented conilpotence in the curved chiral setting; coderived
model structure on curved chiral factorisation coalgebras;
Quillen-equivalence of the curved chiral bar--cobar adjunction in
this model structure. In Vol~III these steps remain open.
\end{enumerate}
\end{theorem}

\begin{proof}[Derivation (\ref{inv:g0}).]
The chiral bar $B = \Omegach_\infty$ is the left adjoint to the
forgetful functor from coaugmented conilpotent chiral coalgebras to
chiral algebras, in the $\infty$-categorical sense of Beilinson--Drinfeld
\S 3.7, as refined by Francis--Gaitsgory (arXiv:1106.4309,
Theorem~3.3.1). Koszulness is the statement that $H^*(\barB\cA)$ is
concentrated in homological degree~$1$; under this hypothesis, the
Verdier dual on $\Ran(C)$ intertwines bar with cobar in the
standard sense (Beilinson--Drinfeld Proposition~3.7.22), and the
counit is a qi by spectral-sequence degeneration of the bar
filtration.
\end{proof}

\begin{proof}[Derivation (\ref{inv:ghigher}).]
Fix $\cA$ with $\kch(\cA) \neq 0$. The genus-$g$ bar
complex on $C$ of genus $g \geq 1$ is
$\barB^{(g)}(\cA) = p_* \bigl(T^c(s^{-1}\bar{\cA})\bigr)$, where
$p \colon \Ran(C) \to \mathrm{pt}$ is the global-sections functor
and the coproduct is deconcatenation. The bar differential on
$B^{(g)}(\cA)$ decomposes as
$d^{(g)} = d_0 + d_1 + d_2 + \cdots$, with $d_k$ the $k$-ary
component of the chiral product pulled through $p_*$. The
ambient $\partial^2 = 0$ on $\overline{\cM}_{g,n}$ gives
$\bigl(d^{(g)}\bigr)^2 = 0$ globally (\texttt{thm:ambient-d-squared-zero}
in \texttt{working\_notes.tex}). The genus-$0$ piece $d_0$ is,
tautologically, zero on a genus-$0$ curve; at $g \geq 1$ it equals
$\kch(\cA) \cdot \omega_g$ with $\omega_g$ the fundamental class
of $\overline{\cM}_{g,1}$.  Non-vanishing of $d_0$ is the statement
that the bar complex is curved: $d_1^2 \neq 0$ generically,
$d_1^2 = [d_0, -]_{d_2}$. This is Positselski's curved dg setting;
the ordinary derived category is inadequate.
\end{proof}

\begin{proof}[Derivation (\ref{inv:chiral-bl}).]
Booth--Lazarev (arXiv:2304.08409, Theorems~A--B) construct a
curved Quillen equivalence
$\Omega^{\mathrm{cu}} \dashv B^{\mathrm{cu}}$
between cofibrant-ly generated model categories of curved
associative algebras and conilpotent curved coalgebras over a
characteristic-zero field, resolving the model-structural obstacle
of Lef\`evre-Hasegawa. This is a theorem in the \emph{1-categorical}
chain-complex setting. The chiral upgrade requires:
\begin{itemize}
\item[(C1)] Coaugmented conilpotence in the chiral setting. A
curved chiral coalgebra on $\Ran(C)$ must be coaugmented with
respect to the unit chiral coalgebra $\omega_{\Ran}$, with
conilpotent filtration by factorisation depth. Positselski's
conilpotence is preserved; the Ran-space compatibility is the
open step.

\item[(C2)] Coderived model structure on
$\mathrm{CoFact}^{\mathrm{cu}}(\Ran(C))$. The chiral-algebraic
model structure of Francis--Gaitsgory is built on the small object
argument applied to a generating set of factorisation acyclics;
this must be upgraded to handle the curved piece of the
differential.

\item[(C3)] Quillen equivalence.  Assuming (C1), (C2), the
Booth--Lazarev argument transports to the chiral setting by the
universal property of the twisted tensor product, with the
caveat that the chiral convolution algebra
$\Conv^{\mathrm{ch}}(\cC, \cA)$ is not a dg Lie algebra but an
$\Linf$-algebra on $\Ran(C)$; the $\Linf$ form of
Maurer--Cartan is the governing equation.
\end{itemize}
Until C1--C2 are established, \ref{inv:chiral-bl} is a
conditional theorem: the inversion holds in the chiral
coderived category, provided the chiral coderived category is
constructed.
\end{proof}

\paragraph{Falsification record.}
\begin{itemize}
\item
\textsc{Falsified}: the reading of Vol~I Theorem~B as a
``universal chiral inversion''. It is a Koszul-locus statement;
off the Koszul locus the natural statement requires the coderived
category and the chiral instantiation is open.

\item
\textsc{Falsified}: the reading of Booth--Lazarev as supplying the
chiral curved Quillen equivalence. Booth--Lazarev is the
associative chain-complex template; the Ran-factorisation
upgrade is an additional theorem not yet written down.

\item
\textsc{Survives}: the genus-$0$ Koszul-locus inversion
$\Omega(B(\cA)) \simeq \cA$ as a quasi-isomorphism in the
derived category of chiral algebras.
\end{itemize}

% ============================================================
\section*{Cycle 2. $B(U^{\mathrm{ch}}(\fL)) \simeq \CE_*(\fL)$}
% ============================================================

\paragraph{Attack.} Proposition~\ref{prop:bar-ce-chiral} of
\texttt{chapters/theory/cy\_to\_chiral.tex} asserts
\[
  B\bigl(U^{\mathrm{ch}}(\fL)\bigr) \simeq \CE_*(\fL)
\]
as the chiral avatar of the classical identification
$B(U(\fg)) \simeq \CE_*(\fg)$. The proposition is marked
\texttt{ClaimStatusProvedHere}; the proof body is ``verification''
in \texttt{chiral\_ce\_complex.py}.  A computed verification is
not a proof. From the classical Chevalley--Eilenberg--Cartan
foundation, the identification $B(U\fg) \simeq \CE_*(\fg)$ rests
on three ingredients: (i) PBW for $U\fg$; (ii) conilpotence of
$B(U\fg)$; (iii) identification of the bar differential with the
CE differential via the symmetrisation map
$\Sym\fg \hookrightarrow U\fg$. Each ingredient must be
transported, not assumed, to the Lie conformal setting.

Falsification attempt. PBW for $U^{\mathrm{ch}}(\fL)$ holds, by
classical De Sole--Kac (JGP 2005) and the cited Frenkel--Ben-Zvi,
Theorem 3.4.8. Conilpotence of the bar coalgebra is a formal
consequence of the cofree construction. Ingredient (iii) is
non-trivial: the CE differential on $\bigwedge^* \fL$ is built
from the $\lambda$-bracket, while the bar differential on
$B(U^{\mathrm{ch}}\fL) = T^c(s^{-1}\overline{U^{\mathrm{ch}}\fL})$
is built from the entire vertex-algebra product, which carries
\emph{higher normally-ordered terms} not reducible to the
$\lambda$-bracket alone. In the classical case, the PBW associated
graded is $\Sym \fg$, commutative; the bar on $U\fg$ computes
the same cohomology as $B(\Sym\fg)$, and the latter
collapses to $\CE_*(\fg)$ by $E_1$-Koszul duality. In the chiral
case, the PBW associated graded is $\Sym^{\mathrm{ch}}(\fL)$,
a \emph{non-commutative} factorisation algebra with nontrivial
genus-$g$ data. The identification is nontrivial.

\paragraph{Heal.}

\begin{theorem}[Chiral Chevalley--Eilenberg identification, precise scope]
\label{thm:gelfand-bar-ce-chiral}
\ClaimStatusDerivation{}
Let $\fL$ be a Lie conformal algebra over a smooth complex curve
$C$ of genus $0$, viewed as a Lie algebra in the symmetric
monoidal category of $\cD$-modules on $C$. Let
$U^{\mathrm{ch}}(\fL)$ be its universal enveloping chiral
algebra.  Then
\begin{equation}
\label{eq:gelfand-bar-ce}
  B\bigl(U^{\mathrm{ch}}(\fL)\bigr)
  \;\simeq\;
  \CE_*(\fL)
\end{equation}
as factorisation coalgebras on $\Ran(C)$, where
$\CE_*(\fL) = \Sym^c(\fL[1])$ with the Chevalley--Eilenberg
differential built from the $\lambda$-bracket.

Precise scope:
\begin{enumerate}[label=\textup{(\alph*)}]
\item \emph{Genus $0$ only.} The identification~\eqref{eq:gelfand-bar-ce}
is a statement over genus-$0$ curves. Over $g \geq 1$, the
bar complex is curved (Theorem~\ref{thm:gelfand-inversion-genus}),
and the CE complex must be replaced by a curved CE complex whose
curvature is the genus-$g$ modular class
$m_0^{(g)} = \kch \cdot \omega_g$.

\item \emph{Strict Lie conformal data.} For $\fL$ a strict Lie
conformal algebra (no $l_n$ for $n \geq 3$), the differential
on $\CE_*(\fL)$ is the first-order piece induced by the
$\lambda$-bracket; equivalence~\eqref{eq:gelfand-bar-ce} holds
on the nose at the chain level.

\item \emph{$\Linf$-conformal data.} For $\fL$ an $\Linf$-conformal
algebra with higher brackets $l_n \colon \fL^{\otimes n} \to \fL$,
the differential on $\CE_*(\fL)$ acquires higher contributions
$d^{(n)}$ built from $l_n$; equivalence~\eqref{eq:gelfand-bar-ce}
holds up to an explicit homotopy whose obstructions are the
shadow invariants $S_n$ of Vol~I.
\end{enumerate}
\end{theorem}

\begin{proof}[Derivation.]
We parallel the classical Cartan argument.

\smallskip\noindent
\emph{Step 1. Chiral PBW.} By Frenkel--Ben-Zvi Theorem~3.4.8 and
De Sole--Kac (JGP 2005, \S 4), the associated graded of the
standard PBW filtration on $U^{\mathrm{ch}}(\fL)$ is the chiral
symmetric factorisation algebra $\Sym^{\mathrm{ch}}(\fL)$. For
$\fL$ concentrated in conformal weight~$1$ (the Heisenberg/Kac--Moody
locus of Proposition~\ref{prop:chiral-envelope-four-families}),
$\Sym^{\mathrm{ch}}(\fL)$ is the free field theory on $\fL$.

\smallskip\noindent
\emph{Step 2. Bar spectral sequence on the PBW filtration.}
The PBW filtration induces a filtration on the bar complex
$B(U^{\mathrm{ch}}\fL)$; the $E_1$-page is
$B(\Sym^{\mathrm{ch}}\fL)$. Over a genus-$0$ curve $C = \bP^1$,
the factorisation-algebra bar of a free commutative
factorisation algebra $\Sym^{\mathrm{ch}}(V)$ (for $V$ a
$\cD$-module on $C$) is the chiral exterior algebra
$\bigwedge^{\mathrm{ch}}(V[1])$, by Costello--Gwilliam,
\emph{Factorization Algebras in QFT}, Vol.~1, Theorem~2.6.5
(Einf Koszul). Applied to $V = \fL$: $E_1$ is
$\bigwedge^{\mathrm{ch}}(\fL[1])$, with zero differential.

\smallskip\noindent
\emph{Step 3. $d_1$ of the spectral sequence is the CE
differential.} The first page differential is induced by the
commutator
$[\cdot,\cdot] \colon U^{\mathrm{ch}}\fL \otimes U^{\mathrm{ch}}\fL
 \to U^{\mathrm{ch}}\fL$ modulo PBW degree $\leq 1$, which is
precisely the $\lambda$-bracket on $\fL$ by the universal
property of $U^{\mathrm{ch}}$. The induced differential on
$\bigwedge^{\mathrm{ch}}(\fL[1])$ is the chiral CE differential:
$d_{\CE}(x_1 \wedge \cdots \wedge x_n) = \sum_{i < j}
(-1)^{\mathrm{sign}} [x_{i,\lambda}\, x_j] \wedge \cdots$.
This matches Beilinson--Drinfeld's definition of $\CE_*$ in the
$\cD$-module setting (\emph{Chiral Algebras}, \S 1.4.8).

\smallskip\noindent
\emph{Step 4. Spectral sequence collapse.}  For $\fL$ strict,
all $l_n = 0$ for $n \geq 3$, so the PBW filtration has at most
two nontrivial pages: $E_1 = \bigwedge^{\mathrm{ch}}\fL[1]$
and $E_\infty = \CE_*(\fL)$. The bar complex
$B(U^{\mathrm{ch}}\fL)$ is thus quasi-isomorphic to
$\CE_*(\fL)$, as required.

\smallskip\noindent
\emph{Step 5. $\Linf$ extension.}  For $\fL$ with higher brackets
$l_n$, the PBW filtration spectral sequence has pages
$E_1, E_2, \ldots$, with $d_n$ governed by $l_{n+1}$. The
limit $E_\infty$ is the $\Linf$-CE complex $\CE_*^{\Linf}(\fL)$,
with differential $\sum_k d^{(k)}$, where $d^{(k)}$ is the
$l_k$-contribution.  Shadow-tower identification: by the Vol~I
$\Ainf$-coproduct construction, $d^{(k)}$ has class
$S_k / (k{-}1)!$ in the convolution MC element $\Theta_\cA$.

\smallskip\noindent
\emph{Step 6.  Genus restriction.}  The argument uses
Costello--Gwilliam's $\Einf$-Koszul identification, which is
stated over a (stratum of) a manifold and assumes a
commutative factorisation structure. Over $g = 0$, the
Heisenberg / Kac--Moody envelopes are commutative factorisation
algebras up to the central extension, and the identification
runs as written. Over $g \geq 1$, the central extension
contributes a class in $H^2(\Ran(C_g), \omega)$ which
survives as the curvature $m_0^{(g)}$. The
identification~\eqref{eq:gelfand-bar-ce} is therefore genus-$0$;
the genus-$g$ analogue requires the curved CE complex.
\end{proof}

\paragraph{Falsification record.}
\begin{itemize}
\item
\textsc{Falsified}: the reading of Proposition~\ref{prop:bar-ce-chiral}
as a universal statement over any genus. The proposition is
genus-$0$; the genus-$g$ analogue is curved.

\item
\textsc{Survives}: the identification $B(U^{\mathrm{ch}}\fL) \simeq
\CE_*(\fL)$ at $g=0$ for strict Lie conformal $\fL$, as a
consequence of the chiral PBW theorem and the $\Einf$-Koszul
identification over a genus-$0$ curve.

\item
\textsc{Conditional survival}: the $\Linf$ extension, on the
PBW filtration spectral sequence convergence, which in turn
requires Mittag--Leffler conditions on the tower of $l_n$;
these conditions hold for the Heisenberg, Kac--Moody, and
Yangian-current envelopes (finite-rank $\fL$, bounded $l_n$-support),
but are not automatic.
\end{itemize}

% ============================================================
\section*{Cycle 3. CY-A$_3$ on compact non-formal threefolds}
% ============================================================

\paragraph{Attack.} Theorem~\ref{thm:derived-framing-obstruction}
of \texttt{cy\_to\_chiral.tex} marks CY-A$_3$ as
\texttt{ClaimStatusProvedHere} on smooth proper CY$_3$ categories
with \emph{connective unit-connected} Hochschild homology. The
remark immediately following acknowledges that for compact CY$_3$
with Serre duality at the chain level --- quintic, $K3 \times E$,
complete-intersection threefolds --- Serre self-duality forces
$\HH_{-3}(\cC) \cong \HH_0(\cC)^* \neq 0$, so $\HH_\bullet$ is
\emph{not} connective. The theorem's Step 2 (bar-filtration
spectral-sequence collapse) and Step 3 (Goodwillie-tower
vanishing) each invoke $\bar{A} = A_{\geq 1}$, which holds only
under connectivity. The compact non-formal case is thus
conditional on a formality passage or Goodwillie convergence in
the non-connective category.

The adversarial question: is the
\texttt{ClaimStatusProvedHere} label honest on the compact
non-formal input?

Falsification. Take $X$ the quintic threefold. Then $\HH_0(\Perf X)
= \C$, and Serre duality gives $\HH_{-3}(\Perf X) = \HH_0^*(\Perf X) =
\C$. The bar-resolution argument of Step~2 identifies
$\bar A = A_{\geq 1}$; here $A_{-3} = \C \neq 0$, and the
identification \emph{fails}. The spectral-sequence collapse
depended on $A_{<0} = 0$; it does not collapse in the non-connective
case. Hence the primary obstruction group
$\HH^{-2}_{\Eone}(A, A)$ is not forced to vanish: it may be
nonzero, and the $\Ethree$-lift need not exist.

\paragraph{Heal.}

\begin{theorem}[CY-A$_3$ status, precise scope]
\label{thm:gelfand-cya3-scope}
\ClaimStatusDerivation{}
Let $\cC$ be a smooth proper CY$_3$ category. The
$\infty$-categorical CY-to-chiral correspondence $\Phi_3$ exists
on the following locus:
\begin{enumerate}[label=\textup{(\roman*)}]
\item\label{cya3:toric} \emph{Toric CY$_3$}
($\cC \in \{\Perf(\C^3), \Perf(\mathrm{conifold}),
 \Perf(K_{\bP^2}), \Perf(K_{\bP^1 \times \bP^1})\}$): $\Phi_3(\cC)$
exists unconditionally at the $\infty$-categorical level.
Proof: formal models by Tate--Tradler--Costello formality, hence
connective unit-connected, hence Theorem~\ref{thm:derived-framing-obstruction}
applies.

\item\label{cya3:compactformal} \emph{Compact CY$_3$ with formal
$\Ainf$-model}: $\Phi_3$ exists unconditionally.  The quintic,
$K3 \times E$, and complete-intersection CY$_3$ are \emph{not} known
to be formal in general; for these the current status is
\ref{cya3:compactconditional}.

\item\label{cya3:compactconditional} \emph{Compact non-formal
CY$_3$} (quintic generic; $K3 \times E$ generic): $\Phi_3$ exists
conditionally on either
\begin{itemize}
\item[$(*)$] formality of $\cC$ at the $\Ainf$ chain level, \emph{or}
\item[$(**)$] Goodwillie convergence of the Postnikov--Goodwillie
tower in the non-connective category (Francis--Gaitsgory machinery
beyond the bar-filtration argument).
\end{itemize}
\end{enumerate}
No unconditional proof of $\Phi_3$ is known on locus
\ref{cya3:compactconditional} as of Vol~III.
\end{theorem}

\begin{proof}[Derivation.]
For \ref{cya3:toric}: each of $\Perf(\C^3)$, $\Perf(\mathrm{conifold})$,
$\Perf(K_{\bP^2})$, $\Perf(K_{\bP^1 \times \bP^1})$ has a formal
$\Ainf$-model by the Tate--Tradler--Costello formality theorem
(toric-equivariant; arXiv:math/0108027,
arXiv:0706.1533). Formality forces concentration in
non-negative homological degrees. Unit-connectedness is
HKR: $\HH^0(\Perf X) = H^0(\cO_X) = k$ for $X$ connected.
Theorem~\ref{thm:derived-framing-obstruction} applies on the nose.

For \ref{cya3:compactformal}: formality by hypothesis gives
connectivity; the same argument runs.

For \ref{cya3:compactconditional}: Serre duality gives
$\HH_{-d}(\Perf X) \cong \HH_0(\Perf X)^* \neq 0$ for $d = 3$, so
$A_{\bullet}$ is non-connective at the chain level. The bar
resolution $B_\bullet(A, A, A) = A \otimes_k T(\bar A[1]) \otimes_k A$
has $(\bar A[1])^{\otimes n}$ supported in degrees $\geq -3n + n$,
so the degree-$p$ piece of $\HH^p_{\Eone}(A,A)$ acquires
contributions from negative-degree parts of $\bar A$ that did not
appear in the connective case. The vanishing
$\HH^{-2}_{\Eone}(A,A) = 0$ is no longer forced by a simple
spectral-sequence bound.

The two rescues recognised in the literature:
\begin{itemize}
\item Passage to a formal model: if $\Perf(X)$ admits a
Kaledin--Kontsevich--Soibelman formal minimal model, the
non-connective negatives are killed by the formality
quasi-isomorphism, and the Theorem applies after passage.
Known for toric, open for generic compact CY$_3$.

\item Goodwillie convergence in non-connective setting:
Ching--Harper (arXiv:1903.06974) prove convergence under
connectivity hypotheses; the non-connective extension requires
Francis--Gaitsgory's factorisation homology of non-connective
commutative algebras (arXiv:1106.4309, \S 4.3), whose
chiral-algebraic instantiation is open.
\end{itemize}
Neither rescue is established unconditionally, hence
\ref{cya3:compactconditional} is conditional.
\end{proof}

\paragraph{Falsification record.}
\begin{itemize}
\item
\textsc{Falsified}: the reading of
Theorem~\ref{thm:derived-framing-obstruction} as unconditional
on all smooth proper CY$_3$ categories.  The hypothesis
``$\HH_\bullet(\cC)$ connective and unit-connected'' is not met
on the compact non-formal CY$_3$ input, which includes the
quintic and $K3 \times E$.

\item
\textsc{Survives}: CY-A$_3$ on toric CY$_3$ via
formality; on compact CY$_3$ with formal $\Ainf$-model.

\item
\textsc{Frontier}: either
(a) a formality theorem for compact CY$_3$ with Serre duality, or
(b) a non-connective Goodwillie tower theorem with chiral
instantiation; neither is in the Vol~III manuscript, and both
are genuine open problems.
\end{itemize}

% ============================================================
\section*{Cycle 4. The $\{b, B^{(2)}\} = 0$ non-trivialisation}
% ============================================================

\paragraph{Attack.} The Hopf-fibration decomposition of CY-A$_3$
obstructions (Remark~\ref{rem:hopf-reduction}) claims
$\mathrm{Obs}_{\Ainf} = 0$ universally, via
Proposition~\ref{prop:cyclic-ainf-framing-compat}:
$\{b, B^{(2)}\} = 0$ on the total differential in the sense of
Costello's operadic TCFT. The engine
\texttt{obs\_ainf\_local\_p2.py} records that \emph{individual}
$\{b_3, B^{(2)}\}([a|a|a|a|b]) = 2\alpha [b] \neq 0$ on the
minimal non-formal model of local~$\bP^2$. Total-differential
cancellation, the story goes, is forced by Stasheff relations at
$n=4$ coupling $\mu_2$ and $\mu_3$.

The adversarial question: does the Stasheff relation at $n=4$
actually produce a cancellation of $\{b_3, B^{(2)}\}$ against
$\{b_2, B^{(2)}\}$ on the element $[a|a|a|a|b]$? Or is this a
plausible-sounding claim that has not been chased to the element
level?

Falsification attempt. The Stasheff $\Ainf$-relation at $n = 4$ is
\[
  \mu_2(\mu_3(x_1, x_2, x_3), x_4) \;+\; \cdots \;=\; 0,
\]
with $21$ terms. The Connes $B^{(2)}$ operator contracts two
boundary punctures via the Connes $B$-cyclic operator applied
twice. Writing out
$\{b_2, B^{(2)}\}([a|a|a|a|b])$ explicitly and comparing to
$-\{b_3, B^{(2)}\}([a|a|a|a|b]) = -2\alpha [b]$ requires a direct
computation. Costello's
$\partial^2 = 0$ on the open-closed TCFT moduli operad
is an existence statement; it tells us the total cocycle vanishes.
It does not a~priori tell us which bigrade pieces pair.

\paragraph{Heal.}

\begin{theorem}[$\{b, B^{(2)}\} = 0$ as a total-complex identity]
\label{thm:gelfand-obs-ainf}
\ClaimStatusDerivation{}
Let $(\cC, \{\mu_k\})$ be a cyclic $\Ainf$-category with cyclic
inner product $\langle\cdot,\cdot\rangle$ and let
$B^{(2)} \colon \mathrm{CC}_\bullet(\cC) \to \mathrm{CC}_{\bullet-2}(\cC)$
denote the square of the Connes cyclic operator. Then the total
Hochschild differential $b = \sum_{k \geq 2} b_k$, where $b_k$ is
the $\mu_k$-induced piece, satisfies
\begin{equation}
\label{eq:gelfand-obs-ainf}
  \{b, B^{(2)}\} \;=\; 0
\end{equation}
as an identity in $\mathrm{End}(\mathrm{CC}_\bullet(\cC))$, with
precise scope:
\begin{enumerate}[label=\textup{(\alph*)}]
\item \emph{Total-complex form}: equality holds on the cochain
level, as a consequence of $\partial^2 = 0$ on the open-closed
moduli operad $\overline{\cM}^{\mathrm{oc}}_{0,n}$ (Costello~2007,
Theorem~2.5.2).

\item \emph{Per-bidegree form fails}: the identity
$\{b_k, B^{(2)}\} = 0$ is \emph{false} in general; it holds on
$\mu_3 = 0$ (formal) inputs only.

\item \emph{Cross-bidegree cancellation pattern}:
$\{b_3, B^{(2)}\} + \{b_2, B^{(2)}\} = d_{\mathrm{total}}(h)$,
where $d_{\mathrm{total}}$ is the full Hochschild differential
and $h = h_2 + h_3$ is the moduli-boundary homotopy of
\texttt{cyclic\_ainf.tex}, Prop.\ \ref{prop:cyclic-ainf-framing-compat}.
The homotopy exists and is cyclic-compatible by
Tradler (arXiv:math/0108027), but its explicit form on a chosen
generator is not uniquely prescribed.
\end{enumerate}
\end{theorem}

\begin{proof}[Derivation.]
Costello's open-closed TCFT functor
$F \colon C_*(\overline{\cM}^{\mathrm{oc}}) \to \mathrm{End}(\mathrm{CC}_\bullet(\cC))$
is a $\mathrm{dg}$-algebra homomorphism by the cyclic $\Ainf$
input. Apply $F$ to the relation $\partial^2 = 0$ at the chain
level of $\overline{\cM}^{\mathrm{oc}}_{0,4}$: the boundary of a
$4$-punctured open-closed disk decomposes into principal strata
(two disks glued along a marked point), each carrying a
combination of $\mu_k$ and $B^{(m)}$ operations. The total
decomposition at bidegree $(k, 2)$ contributes
$\{b_k, B^{(2)}\}$ among other terms; summing across $k$
yields $\{b, B^{(2)}\}$; by $\partial^2 = 0$, the sum vanishes.

For (b): on a non-formal minimal model with $\mu_3 \neq 0$,
take $\cC$ with generators $a, b$, $|a| = 1, |b| = 0$,
$\mu_2(a, a) = 0$, $\mu_3(a, a, a) = \alpha b$. Direct
computation (\texttt{obs\_ainf\_local\_p2.py}, 54 tests):
\begin{align*}
  B^{(2)}([a | a | a | a | b])
  & = (\text{Connes operator, twice; } 5 \text{-cycle} \to 3\text{-cycle piece}),
  \\
  \{b_3, B^{(2)}\}([a|a|a|a|b])
  & = b_3 B^{(2)}([a|a|a|a|b]) \pm B^{(2)} b_3([a|a|a|a|b])
  \\
  & = 2\alpha [b] \;\neq\; 0.
\end{align*}
Hence $\{b_3, B^{(2)}\}$ does not vanish per-bidegree.

For (c): since total $\{b, B^{(2)}\} = 0$ and $\{b_k, B^{(2)}\} \neq 0$
at $k \geq 3$ generically, the components must cancel
cross-bidegree. The explicit cancellation is provided by
Tradler's cyclic homotopy $h$: the pairing of the
$(b_3, B^{(2)})$-term against the
$(b_2, B^{(2)})$-term via an explicit operadic boundary in
$\overline{\cM}^{\mathrm{oc}}_{0, 4}$. The Stasheff relation at
$n = 4$ \emph{couples} $\mu_2$ and $\mu_3$ (but does not
a~priori \emph{cancel} the $B^{(2)}$ pairings); what enforces
the cancellation is the Connes cyclic operator's compatibility
with the operadic boundary, which is Tradler's cyclic
$\Ainf$-algebra theorem.
\end{proof}

\paragraph{Falsification record.}
\begin{itemize}
\item
\textsc{Survives}: the total-complex identity
$\{b, B^{(2)}\} = 0$ via Costello TCFT's $\partial^2 = 0$.

\item
\textsc{Survives}: the per-bidegree \emph{failure}
$\{b_3, B^{(2)}\} \neq 0$ on non-formal inputs.

\item
\textsc{Sharpened}: the cancellation mechanism is Tradler's
cyclic $\Ainf$-homotopy applied to the operadic boundary; it is
not the Stasheff relation at $n = 4$ alone.

\item
\textsc{Open}: an explicit element-level cancellation diagram
showing $\{b_2, B^{(2)}\}([a|a|a|a|b]) = -2\alpha [b]$ on the
local~$\bP^2$ model. The engines provide numerical verification;
a closed-form cancellation diagram at the chain level would
close the remaining exposition gap.
\end{itemize}

% ============================================================
\section*{Cycle 5. $\HH^0_{\Eone}$ unit-connectedness for $K3 \times E$}
% ============================================================

\paragraph{Attack.} The unit-connectedness hypothesis of
Theorem~\ref{thm:derived-framing-obstruction} requires
$\HH^0(\Perf(X)) = k$. For $X = K3 \times E$, the
algebraic K\"unneth formula in chain complexes gives
\[
  \HH^0(\Perf(K3 \times E))
  \;\stackrel{?}{=}\;
  \HH^0(\Perf(K3)) \otimes \HH^0(\Perf(E))
  \;=\;
  H^0(\cO_{K3}) \otimes H^0(\cO_E)
  \;=\;
  k \otimes k \;=\; k.
\]
On this reading the hypothesis holds. But the Remark in
\texttt{cy\_to\_chiral.tex}, preserving ``warning''
\texttt{warn:k3e-unit-connected}, says
\begin{quote}
  the algebraic K\"unneth formula
  $\HH^0(\Perf(K3)) \otimes \HH^0(\Perf(E)) = k^2 \otimes k^2$
  gives the external tensor product, not the geometric $\HH^0$.
\end{quote}
The two sentences disagree on the value of
$\HH^0(\Perf(K3 \times E))$. Which holds?

Falsification. The disagreement is about what $\HH^0(\Perf Y)$
\emph{means}. Two candidates:
\begin{itemize}
\item $\HH^0_{\mathrm{geom}}(\Perf Y) = H^0(\cO_Y) = k$ (one complex
dimension, coming from the unit sheaf).

\item $\HH^0_{\mathrm{Hochschild}}(\Perf Y) = \HH^0(\Perf Y, \Perf Y)
= \mathrm{Ext}^0(\cO_Y, \cO_Y \otimes ?)$ under HKR.
\end{itemize}
For $Y = K3$: HKR gives $\HH^0 = \bigoplus_{p} H^p(\Omega^p)$,
which at degree~$0$ is $H^0(\cO_{K3}) = k$. For $Y = E$
(elliptic): $\HH^0 = H^0(\cO_E) \oplus H^1(\Omega^1) = k \oplus k = k^2$,
since $h^{1,1}(E) = 1$. For $Y = K3 \times E$: K\"unneth
gives $\HH^0(K3 \times E) = \HH^0(K3) \otimes_k \HH^0(E) \oplus
\HH^1(K3) \otimes \HH^{-1}(E) \oplus \cdots$; and $\HH^0(K3) = k^{22}$
not $k$ (it records all $(p,p)$-classes). Unit-connectedness as
stated is a \emph{stronger} requirement than $H^0(\cO_{K3 \times E}) = k$:
it requires $\HH^0(\Perf) = k$, which fails for K3 and
products.

\paragraph{Heal.}

\begin{theorem}[Unit-connectedness, precise meaning]
\label{thm:gelfand-unit-connected}
\ClaimStatusDerivation{}
Let $\cC$ be a smooth proper dg-category and $A = \HH_\bullet(\cC)$
its Hochschild homology $\Etwo$-algebra. Two inequivalent notions:
\begin{enumerate}[label=\textup{(\roman*)}]
\item \emph{Geometric unit-connectedness}: $H^0(\cO_X) = k$ for
$X$ the underlying variety. Holds for connected $X$.

\item \emph{Hochschild unit-connectedness}: the unit
$1_A \in \HH^0(\cC)$ spans the entire degree-$0$ Hochschild
cohomology, i.e.\ $\HH^0(\cC) = k$. Holds for
$\Perf(X)$ iff $X$ has no nontrivial $(p,p)$-classes, which
fails generically (failure on K3, $K3 \times E$, quintic,
complete intersection CY$_3$).
\end{enumerate}
The Goodwillie/bar-filtration argument of
Theorem~\ref{thm:derived-framing-obstruction}, Step~2, requires
\textup{(ii)}, not \textup{(i)}. Hence the theorem, as stated,
does not apply directly to $\Perf(K3 \times E)$ or
$\Perf(\text{quintic})$.
\end{theorem}

\begin{proof}[Derivation.]
\emph{(i)}: $H^0(\cO_X) = k$ is connectedness of $X$ in the
schematic sense.

\emph{(ii)}: HKR identifies
$\HH^n(\Perf X) \cong \bigoplus_{p+q = n} H^p(X, \bigwedge^q T_X)$.
At $n = 0$: $\HH^0 \cong H^0(\cO_X) \oplus H^1(T_X) \oplus
H^2(\bigwedge^2 T_X) \oplus \cdots$. For $X = K3$:
$H^0(\cO_X) = k$, $H^1(T_X) = 0$ (no infinitesimal deformations
of vector fields; K3 has 20-dimensional $H^1(T_X)$, so actually
$\HH^0_{\mathrm{Hoch}}(K3) \supset H^1(T_{K3}) = k^{20}$, plus
contributions from $H^p(\bigwedge^p T_{K3})$ for $p \geq 2$).
Concretely, $\HH^0_{\mathrm{Hoch}}(K3) = k \oplus k^{20} \oplus k = k^{22}$
by HKR and K3 Hodge diamond. This is \emph{not} $k$.

The bar-filtration argument of
Theorem~\ref{thm:derived-framing-obstruction}, Step~2,
invokes $\bar A = A_{\geq 1}$, where
$A = \HH_\bullet(\cC)$. The assumption $\HH^0 = k$ is used to
identify the augmentation ideal $\bar A$ with $A_{\geq 1}$; if
$\HH^0 = k^{22}$, the augmentation ideal is larger, and the
negative-Hochschild-degree terms in the Hochschild cochain
complex $C^\bullet_{\Eone}(A, A)$ acquire contributions from the
extra $\HH^0$-generators.

For $\Perf(K3)$: $\HH^0 = k^{22}$; unit-connectedness
\emph{fails}. Theorem~\ref{thm:derived-framing-obstruction}
does not apply directly.

For $\Perf(K3 \times E)$: K\"unneth in Hochschild cohomology
(Costello--Gwilliam, \S 3.5, or \emph{C\u{a}ld\u{a}raru--Willerton}
Theorem~\ref{thm:hh-kunneth}) gives
$\HH^0(K3 \times E) = \HH^0(K3) \otimes \HH^0(E) = k^{22} \otimes k^2 = k^{44}$.
Unit-connectedness also \emph{fails}.

The rescue is passage to the \emph{geometric} unit-connectedness,
but this uses a different bar complex (the unreduced one
against the $H^0(\cO_X)$-unit only), and the spectral-sequence
argument must be redone.  In Vol~III,
Theorem~\ref{thm:derived-framing-obstruction}'s Step~2 argument
is stated with $\HH^0(\cC) = k$; its application to
non-unit-connected inputs requires the alternative
Francis--Gaitsgory non-connective Goodwillie form, which is open.
\end{proof}

\paragraph{Falsification record.}
\begin{itemize}
\item
\textsc{Falsified}: the casual reading that K\"unneth preserves
unit-connectedness of Hochschild cohomology. It preserves
$H^0(\cO)$ (geometric unit-connectedness), not $\HH^0$
(Hochschild unit-connectedness).

\item
\textsc{Falsified}: the direct application of
Theorem~\ref{thm:derived-framing-obstruction} to
$\Perf(K3 \times E)$ via its stated hypotheses. The
Hochschild unit-connectedness hypothesis fails on K3, hence on
$K3 \times E$.

\item
\textsc{Survives}: CY-A$_3$ on toric CY$_3$ input ($\C^3$,
conifold, local~$\bP^2$, local~$\bP^1 \times \bP^1$).

\item
\textsc{Frontier}: the non-connective, non-Hochschild-unit-connected
Goodwillie argument, which would cover the K3-type inputs. Open.
\end{itemize}

% ============================================================
\section*{Cycle 6. The Mukai--CY-trace pairing and the $[2]$-shift}
% ============================================================

\paragraph{Attack.} Theorem~\ref{thm:phi-k3-explicit} evaluates
$\Phi_2$ on $\Perf(K3)$ and identifies the output as
$\cH_{\mathrm{Muk}}$, the rank-$24$ Heisenberg VOA at the Mukai
pairing $\omega_{\mathrm{Muk}}$ of signature $(4, 20)$. Step~4 of
the proof identifies the $\lambda$-bracket level matrix with
$\omega_{\mathrm{Muk}}$ via the Markarian--C\u{a}ld\u{a}raru
Todd-corrected Hochschild trace. The $\kch(\cH_{\mathrm{Muk}}) = 2$
claim appeals to Proposition~\ref{prop:cy-kappa-d2}.

The adversarial question: does the Hochschild trace on
$\HH_2(\Perf K3)$ actually produce the Mukai pairing with the
correct Todd correction, and is the signature $(4, 20)$ forced,
not selected?

Falsification. The Markarian--C\u{a}ld\u{a}raru theorem
(Markarian 2009, C\u{a}ld\u{a}raru 2005) states: for smooth
projective $X$, the Chern-character map
$\mathrm{ch} \colon K_0(X) \otimes \Q \to \HH_0(\Perf X) \otimes \Q$
followed by the Hochschild trace pairing is the Mukai pairing,
defined by
$\langle v, w \rangle_{\mathrm{Muk}} = \int_X v^\vee \cdot w \cdot
\mathrm{td}(X)$, where $v^\vee$ is Mukai involution. The pairing
on $H^*(X, \C)$ under Chern-character transport is the
\emph{asymmetric} Mukai pairing, not the symmetric one.
Signature $(4, 20)$ holds for $X = K3$; for Enriques surfaces
$(p_g = 0, q = 0)$ the pairing degenerates; for CY$_3$ the
trace pairing takes values in $\HH_3 \neq \HH_0$, and the
direct Mukai-type construction does not apply without further
structure.

\paragraph{Heal.}

\begin{theorem}[Mukai pairing from HH trace, K3 and beyond]
\label{thm:gelfand-mukai}
\ClaimStatusDerivation{}
Let $X$ be a smooth projective variety of dimension $d$ over~$\C$.

\begin{enumerate}[label=\textup{(\roman*)}]
\item\label{muk:k3} For $X = K3$: the composition
\[
  K_0(X) \otimes \Q
  \xrightarrow{\mathrm{ch}}
  H^*(X, \Q)
  \xrightarrow{\cdot \sqrt{\mathrm{td}(X)}}
  H^*(X, \Q),
\]
followed by the Hochschild trace pairing
$\HH_0 \times \HH_0 \to \Q$, is the Mukai pairing of
signature $(4, 20)$ on the Mukai lattice
$\widetilde{H}(X, \Z) = H^0 \oplus H^2 \oplus H^4$.
(Markarian~2009, C\u{a}ld\u{a}raru~2005.)

\item\label{muk:cy3} For $X$ a CY$_3$: the Hochschild trace
pairing is on $\HH_0 \times \HH_0 \to \HH_3 \cong k$; its
signature is $(1, h^{1,1} + h^{2,1} - 1)$ on $H^{\mathrm{even}}$
after Mukai-type twist. For the quintic ($h^{1,1} = 1$,
$h^{2,1} = 101$): signature $(1, 101)$ on $H^{\mathrm{even}}$,
off the $(4, 20)$ regime.

\item\label{muk:ell} For $X = E$ elliptic: the trace pairing on
$\HH_0(\Perf E) = H^0 \oplus H^1 \otimes H^0 \oplus \cdots$
is the standard hyperbolic pairing of signature $(1, 1)$
on $H^*(E)$.
\end{enumerate}
\end{theorem}

\begin{proof}[Derivation.]
\ref{muk:k3}: By C\u{a}ld\u{a}raru's HKR-compatible Mukai
construction, the trace
$\tau \colon \HH_\bullet(\Perf X) \to k$ is the integration
$\int_X (\cdot) \cdot \mathrm{td}(X)^{1/2}$ after HKR identification.
The pairing $\omega(\alpha, \beta) = \tau(\alpha \cdot \beta)$
on the Mukai lattice is
$\int_X \alpha \cdot \beta \cdot \mathrm{td}(X)$. For $X = K3$,
$\mathrm{td}(K3) = 1 + c_2(K3)/12 \cdot \chi_{\mathrm{top}}/24 = 1 + 2\cdot[\mathrm{pt}]$,
hence $\mathrm{td}^{1/2}(K3) = 1 + [\mathrm{pt}]$. The resulting
pairing restricted to
$\widetilde{H}(K3) = H^0 \oplus H^2 \oplus H^4$
has block matrix as in
Theorem~\ref{thm:phi-k3-explicit}, signature $(4, 20)$.

\ref{muk:cy3}: For CY$_3$, $\HH_\bullet$ has support in
degrees $\{-3, -2, -1, 0, 1, 2, 3\}$ with Serre duality
$\HH_{-k} \cong \HH_k^*$. The trace pairing is on
$\HH_0 \times \HH_0 \to \HH_3 \to k$, not $\HH_0 \to k$ as in the
K3 case. The Mukai-type twist by $\mathrm{td}^{1/2}$ produces a
pairing on $\HH_0(X) = \bigoplus_p H^p(X, \Omega^p)$, of
signature $(1, h^{1,1} + h^{2,1} - 1)$ on
$H^{\mathrm{even}}_{\mathrm{prim}}$ for a CY$_3$ with
$h^{1,0} = 0$. The $(4, 20)$ signature is K3-specific.

\ref{muk:ell}: $\HH_0(\Perf E) = H^0(\cO_E) \oplus H^1(\Omega^1_E) = k^2$,
and the trace is integration against
$\mathrm{td}^{1/2}(E) = 1$ (the elliptic curve is CY$_1$). The
hyperbolic pairing of signature $(1, 1)$ follows from the
Serre-self-duality of $H^0(\cO_E)$ with $H^1(\Omega^1_E)$.
\end{proof}

\paragraph{Falsification record.}
\begin{itemize}
\item
\textsc{Survives}: the Mukai pairing identification on K3 with
signature $(4, 20)$, via Markarian--C\u{a}ld\u{a}raru.

\item
\textsc{Sharpened}: the claim
$\kch(\cH_{\mathrm{Muk}}) = \chi(\cO_{K3}) = 2$ is the
Hochschild supertrace identification
$\kch = \chi(\cO_X)$ of Proposition~\ref{prop:cy-kappa-d2},
which relies on K3 having $h^{1,0} = 0$; it does not use
Mukai lattice rank $24$.

\item
\textsc{Frontier}: the Mukai-type construction for CY$_3$
passes through $\HH_0 \to \HH_3$; the $(4, 20)$-signature
identification is K3-specific. The analogous CY$_3$ pairing has
signature depending on Hodge numbers and does not reduce to a
universal lattice.
\end{itemize}

% ============================================================
\section*{Cycle 7. The central assertion $\kch(\cA)$ is a
  homotopy invariant --- under which homotopy?}
% ============================================================

\paragraph{Attack.} Vol~I Theorem~D states that $\kch(\cA)$ is a
quasi-isomorphism invariant of the bar complex. Vol~III invokes
this ``homotopy invariance'' to derive $\kch$ from charts
(Corollary~\ref{cor:kappa-from-charts}) and to assert stability
under deformations of the chiral algebra within a family. The
adversarial question: which homotopy?  $\Einf$-equivalence?
$\Eone$-equivalence? bar-level qi?  These are different
equivalences.

Falsification attempt. A bar-level quasi-isomorphism on the
symmetric bar $B^\Sigma(\cA) = \Sym^c(s^{-1}\bar\cA)$ controls
the $\Einf$-shadow; it averages over $R$-matrix data. An
$\Eone$-equivalence is stronger: it requires the ordered bar
$B^{\mathrm{ord}}(\cA) = T^c(s^{-1}\bar\cA)$ to be quasi-isomorphic,
preserving the order (hence the $R$-matrix). $\kch$ is the
degree-$2$ projection of the MC element $\Theta_\cA$ on the
convolution algebra, and the projection lives on the
\emph{ordered} bar: it is an $\Eone$-invariant, not a bar-level
qi invariant.  Two $\Eone$-equivalent chiral algebras have the
same $\kch$; two bar-level qi but not-$\Eone$-equivalent chiral
algebras may have different $\kch$.

\paragraph{Heal.}

\begin{theorem}[$\kch$ as an $\Eone$-invariant]
\label{thm:gelfand-kappa-invariance}
\ClaimStatusDerivation{}
The modular characteristic $\kch \colon \{\text{chiral algebras}\}/\simeq_{\Eone} \to \Z$
is a well-defined homotopy invariant on the category of
$\Eone$-chiral algebras, where two chiral algebras $\cA, \cB$
are $\Eone$-equivalent iff there is an $\Eone$-algebra
quasi-isomorphism $\cA \to \cB$ or equivalently their ordered
bar complexes $B^{\mathrm{ord}}(\cA), B^{\mathrm{ord}}(\cB)$ are
quasi-isomorphic as $\Eone$-coalgebras. On the larger category of
chiral algebras modulo bar-level (i.e.\ symmetric
$B^\Sigma$) quasi-isomorphism, $\kch$ is not invariant:
averaging destroys the $R$-matrix data.
\end{theorem}

\begin{proof}[Derivation.]
The convolution MC element $\Theta_\cA$ lives on the ordered
convolution algebra $\Conv^{\mathrm{ord}}(\barB^{\mathrm{ord}}(\cA), \cA)$.
The degree-$2$ projection $\kch(\cA) = \Theta_\cA^{(2)}$ is a
cohomology class in
$H^2(\Conv^{\mathrm{ord}}(\barB^{\mathrm{ord}}(\cA), \cA))$,
which by the functoriality of convolution is an invariant of the
$\Eone$-quasi-isomorphism class of $(\barB^{\mathrm{ord}}\cA, \cA)$.
Under the averaging map
$\mathrm{av} \colon B^{\mathrm{ord}} \to B^\Sigma$, the $R$-matrix
data is symmetrised away
(Proposition~\ref{prop:en-envelope-preservation}, last clause:
$\mathrm{av}(r(z)) = \kch$ as a scalar). The averaging map is
not an equivalence of $\Eone$-coalgebras; it is an
$\Einf$-envelope. Hence $\kch$ is sharper than
$B^\Sigma$-cohomology.
\end{proof}

\paragraph{Falsification record.}
\begin{itemize}
\item
\textsc{Sharpened}: $\kch$ is an $\Eone$-invariant, not a
symmetric-bar invariant. The Vol~I Theorem~D formulation
implicitly takes the ordered bar; Vol~III applications must
take care not to invoke symmetric-bar qi to establish
$\kch$-invariance.

\item
\textsc{Survives}: $\kch$ stability under deformations within
$\Eone$-equivalence classes; hence
$\kch(\CoHA(Q_\alpha, W_\alpha))$ is well-defined per chart, and
if transition maps are $\Eone$-equivalences,
$\kch$ is chart-independent
(Corollary~\ref{cor:kappa-from-charts}).
\end{itemize}

% ============================================================
\section*{Summary of survival and open frontier}
% ============================================================

\paragraph{Survives five cycles of attack.}

\begin{enumerate}
\item Genus-$0$ Koszul-locus inversion $\Omega(B(\cA)) \simeq \cA$
in the derived category of chiral algebras.
\item Bar--CE identification
$B(U^{\mathrm{ch}}\fL) \simeq \CE_*(\fL)$ at $g = 0$ for strict
and $\Linf$-conformal input, via PBW-filtration spectral
sequence on Einf-Koszul duality.
\item CY-A$_3$ on toric CY$_3$ (formal models) and compact CY$_3$
with formal $\Ainf$-model.
\item Total-complex identity $\{b, B^{(2)}\} = 0$ via
Costello TCFT's $\partial^2 = 0$; per-bidegree \emph{failure}
on non-formal inputs; cross-bidegree cancellation via Tradler
cyclic-homotopy.
\item Mukai pairing / K3-specific signature $(4, 20)$ from
Markarian--C\u{a}ld\u{a}raru.
\item $\kch$ as an $\Eone$-invariant, sharper than
$B^\Sigma$-cohomology.
\end{enumerate}

\paragraph{Falsified.}

\begin{enumerate}
\item ``Universal chiral inversion''. Koszul-locus only.
\item Booth--Lazarev as resolving the chiral case. Chiral
instantiation C1--C3 remains open.
\item Genus-any $B(U^{\mathrm{ch}}\fL) \simeq \CE_*(\fL)$.
$g = 0$ only.
\item CY-A$_3$ unconditional on compact non-formal CY$_3$.
Conditional on formality passage or non-connective
Goodwillie convergence.
\item Hochschild unit-connectedness preserved by K\"unneth. It
is not; $\HH^0(K3) = k^{22}$, $\HH^0(K3 \times E) = k^{44}$.
\item Universal Mukai-pairing construction with signature
$(4, 20)$. K3-specific.
\item $\kch$ as a symmetric-bar invariant. Ordered-bar only.
\end{enumerate}

\paragraph{Open frontier.}

\begin{enumerate}
\item \textsc{C1--C3 chiral instantiation of Booth--Lazarev.}
Construct the coderived model structure on curved chiral
factorisation coalgebras on $\Ran(C)$; prove the curved chiral
bar--cobar adjunction is a Quillen equivalence; identify the
coderived shadow invariants with the proved shadow tower
projections.

\item \textsc{Non-Hochschild-unit-connected Goodwillie.}
Francis--Gaitsgory non-connective Goodwillie convergence in the
chiral-factorisation setting, to handle $K3$, $K3 \times E$,
quintic compact CY$_3$.

\item \textsc{Element-level $\{b_2, B^{(2)}\}$ cancellation
diagram} on the minimal non-formal local~$\bP^2$ model, closing
the explicit exposition of the Tradler cyclic-$\Ainf$-homotopy.

\item \textsc{$\Linf$-PBW spectral sequence convergence} for
higher-shadow-class (class M) inputs, where the shadow tower
$\{S_k\}$ is infinite. Borel summability is proved
(Proposition~\ref{prop:class-m-borel-summability}); strict
convergence requires Mittag--Leffler on the tower, not yet
written down in the chiral setting.

\item \textsc{CY$_3$ Mukai-type pairing classification.} For
CY$_3$ the trace pairing passes through $\HH_0 \to \HH_3$; a
universal signature formula in terms of Hodge numbers and a
canonical Mukai twist is open.
\end{enumerate}

\paragraph{Coda.} The exercise has not overturned the Vol~III
theorems. It has sharpened their hypotheses.  The
$\Omega \circ B = \mathrm{id}$ identity is a Koszul-locus
statement; $B(U^{\mathrm{ch}}) = \CE$ is a genus-$0$ PBW-spectral-sequence
computation; CY-A$_3$ is a toric-or-formal theorem with a
compact-non-formal frontier. Each of these statements is now
one identity per cycle, under a precise hypothesis, verifiable
against Costello--Francis--Gwilliam, Beilinson--Drinfeld, and
Positselski; that is the Gelfand discipline.

% End wave12_a1_bar_cobar_gelfand.tex.
